\documentclass[11pt, a4paper]{article}

\usepackage[T1]{fontenc}
\usepackage[utf8]{inputenc}
\usepackage{tgpagella}
\usepackage[spanish, es-noshorthands]{babel}
\usepackage{amsmath, amssymb, bm}
\usepackage{tabularx}
\usepackage{booktabs}
\usepackage[most]{tcolorbox}
\usepackage[margin=2.0cm]{geometry}
\usepackage{ragged2e}
\usepackage{fancyhdr}

\pagestyle{fancy}
\fancyhf{}
\setlength{\headheight}{16pt}
\lhead{\textbf{UCB Sede Tarija} | Ing. Mecatrónica}
\chead{}
\rhead{IMT-121 Circuitos Electrónicos I | Guía Autónoma W16S2}
\lfoot{Prof: M.Sc. Ing. Bernardo Quiroga Turdera}
\rfoot{Página \thepage}
\renewcommand{\headrulewidth}{0.4pt}

\definecolor{ucbblue}{RGB}{0, 51, 102}

\begin{document}

\begin{tcolorbox}[colback=ucbblue!5!white, colframe=ucbblue, title=\textbf{Artefacto A5a: Guía Autónoma de Trabajo (W16S2 — 18/11/2026)}]
    \centering \textbf{Rutas A (Con GO Presencial) y B (Analítica sin Energización)}
\end{tcolorbox}

\noindent\textit{Objetivo: Completar predicciones analíticas propias y, SOLO si existe un supervisor institucional habilitado que firme el GO, realizar la adquisición de datos de ganancia y fase.}

\section*{Fórmulas Clave}
$$ \tau=RC, \quad f_c=\frac{1}{2\pi RC}, \quad x=\frac{f}{f_c}, \quad |H|=\frac{1}{\sqrt{1+x^2}}, \quad \phi=-\arctan(x) $$
Si logra medir: $\quad |H|_{\text{med}} = \frac{V_{out,pp}}{V_{in,pp}}, \quad \phi_{\text{med}} = -360^\circ \frac{\Delta t}{T}$

\section*{Cronograma de Sesión (90 min)}
\textbf{Minutos 0–10 (Lectura y Autochequeo):} Identifique los componentes R y C de su grupo. Verifique la presencia del supervisor institucional habilitado. Si no hay supervisor o no da GO, asuma \textbf{Ruta B}. \par
\textbf{Minutos 10–30 (Cálculos Individuales):} Calcule $\tau$, $f_c$ y proponga 3 frecuencias ($0.1f_c, f_c, 10f_c$). Calcule $|H|$ y $\phi$ predictivos. Registre en la hoja A2. \par
\textbf{Minutos 30–70 (Ejecución según Ruta):}
\begin{itemize}
    \item \textbf{Ruta A (Con GO):} Con salida en OFF arme el circuito, unifique GNDs, consiga firma GO. Alimente el RC, lea $V_{in,pp}$, $V_{out,pp}$ y $\Delta t$. Guarde captura.
    \item \textbf{Ruta B (NO-GO):} Circuito apagado. Repase las mediciones pasivas (Multímetro) de $R$ y $C$. Elabore sus predicciones analíticas completas indicando el bloqueo técnico.
\end{itemize}
\textbf{Minutos 70–85 (Comparación):} Compare datos con predicciones. Para Ruta B, discuta qué evidencia falta. \par
\textbf{Minutos 85–90 (Cierre y Entrega):} Adjunte hoja A2 y evidencias. Cierre según protocolo de seguridad.

\section*{Checklist de Entrega}
[ ] Integrantes e identificación de ruta. \quad [ ] $R, C$ propios, $\tau, f_c$. \quad [ ] GO y datos reales (Si Ruta A). \par
[ ] Explicación de obstáculo (Si Ruta B). \quad [ ] Pendientes y solicitud de recuperación.

\end{document}
